
\section{Die Dehnungsfamilie der Viererordnungen}

Wir betrachten die Familie
\[
Q_k(p):=(p,p+2,p+6+12k,p+8+12k),
\qquad k\in\mathbb N_0,
\]
sofern alle vier Einträge prim sind. Diese Familie erhält die Zwillingsflügelstruktur sowie die EABC-Klassenarchitektur der primalen Grundform
\[
Q_0(p)=(p,p+2,p+6,p+8).
\]

\subsection{Dehnungsparameter}

Der Parameter
\[
m_k:=6+12k
\]
wird als Flügel- oder Dehnungsparameter interpretiert. Die \(k\)-te Stufe beschreibt eine strukturtreue Deformation derselben Vierergrundordnung. Äquivalent kann man schreiben
\[
Q_k(p)=(p,p+2,p+m_k,p+m_k+2).
\]

\subsection{Stufenübersicht}

Die ersten fünf Stufen lauten:
\[
Q_0(p)=(p,p+2,p+6,p+8),
\]
\[
Q_1(p)=(p,p+2,p+18,p+20),
\]
\[
Q_2(p)=(p,p+2,p+30,p+32),
\]
\[
Q_3(p)=(p,p+2,p+42,p+44),
\]
\[
Q_4(p)=(p,p+2,p+54,p+56).
\]

Dabei ist \(Q_0\) die Urstufe oder primale Grundform, während \(Q_k\) für \(k\ge 1\) als Dehnungsordnungen derselben kompatiblen Viererarchitektur interpretiert werden.

\subsection{Metrischer Kreuzrest}

Für die lineare logarithmische Geometrie gilt in der Kreuzform asymptotisch
\[
\Delta E_{\mathrm{metrisch}}(p,k)
\sim
\frac{64+288k+288k^2}{p^2}.
\]
Insbesondere wächst der metrische Dehnungsrest quadratisch mit dem Dehnungsrang \(k\).

Die ersten Spezialfälle lauten:
\[
\Delta E_{\mathrm{metrisch}}(p,0)\sim \frac{64}{p^2},
\qquad
\Delta E_{\mathrm{metrisch}}(p,1)\sim \frac{640}{p^2},
\]
\[
\Delta E_{\mathrm{metrisch}}(p,2)\sim \frac{1792}{p^2},
\qquad
\Delta E_{\mathrm{metrisch}}(p,3)\sim \frac{3520}{p^2},
\qquad
\Delta E_{\mathrm{metrisch}}(p,4)\sim \frac{5824}{p^2}.
\]

\subsection{EABC-Strukturspannung}

In der EABC-Einbettung bleibt der führende Term der klassengeometrischen Kreuzspannung universell:
\[
\Delta E_{EABC}(p,k)\sim 4(\log p)^2.
\]
Damit ist die grobe Fehlverschaltung der Kreuzform unabhängig vom Dehnungsrang, während \(k\) nur in Untertermen erscheint.

\subsection{Walter-Dehnungssplits}

Für die kanonischen Triangulierungen erhält man asymptotisch
\[
\Delta W_{AC}(p,k)\sim \Delta W_{BE}(p,k)\sim \frac{(6+12k)\log p}{10p}.
\]
Die Walter-Kanäle wachsen somit linear mit dem Dehnungsrang \(k\) und messen die dehnungsinduzierte Besetzungsasymmetrie.

Für die ersten Stufen ergibt sich damit:
\[
k=0:\ \Delta W \sim \frac{3\log p}{5p},
\qquad
k=1:\ \Delta W \sim \frac{9\log p}{5p},
\]
\[
k=2:\ \Delta W \sim \frac{3\log p}{p},
\qquad
k=3:\ \Delta W \sim \frac{21\log p}{5p},
\qquad
k=4:\ \Delta W \sim \frac{27\log p}{5p}.
\]

\subsection{Morley-Dehnungsprofile}

Die Morley-Splits erfüllen für festes \(k\) und \(p\to\infty\)
\[
\Delta M_{AC}(p,k)\to 0,
\qquad
\Delta M_{BE}(p,k)\to 0.
\]
Somit laufen alle kanonischen Dreiecke asymptotisch in dieselbe Morley-Standardgestalt zurück. Die \(k\)-Abhängigkeit erscheint nur in den endlichen Gestaltkorrekturen.

Plausibel ist eine Kleinordnung der Form
\[
\Delta M_{AC}(p,k),\ \Delta M_{BE}(p,k)=O\!\left(\frac{m_k}{p}\right)
\]
oder feiner
\[
O\!\left(\frac{m_k\log p}{p}\right),
\]
was numerisch für allgemeine \(k\) noch gesondert zu testen wäre.

\subsection{Zusammenfassende Skalengesetze}

Für die \(k\)-te Dehnungsordnung gelten damit die Leitgesetze
\[
\Delta E_{\mathrm{metrisch}}(p,k)
\sim
\frac{64+288k+288k^2}{p^2},
\]
\[
\Delta E_{EABC}(p,k)\sim 4(\log p)^2,
\]
\[
\Delta W_{AC}(p,k),\ \Delta W_{BE}(p,k)\sim \frac{(6+12k)\log p}{10p},
\]
\[
\Delta M_{AC}(p,k),\ \Delta M_{BE}(p,k)\to 0
\qquad (p\to\infty,\ k\text{ fest}).
\]

\subsection{Interpretation als Dehnungsfamilie}

Die Familie \(\{Q_k(p)\}_{k\ge 0}\) bildet eine strukturtreue Dehnungsfamilie. Über alle Stufen hinweg bleiben erhalten:
\begin{itemize}
\item die Zwillingsflügel,
\item die EABC-Klassenarchitektur,
\item die lineare ptolemäische Schließung der kompatiblen Form,
\item die kanonischen Triangulierungen,
\item die asymptotische Morley-Standardgestalt.
\end{itemize}

Variabel sind dagegen die dehnungsbedingten Korrekturterme:
\begin{itemize}
\item der metrische Kreuzrest wächst quadratisch in \(k\),
\item die Walter-Besetzung wächst linear in \(k\),
\item die Morley-Feinverzerrungen verschwinden asymptotisch, tragen aber für endliche Größen ein \(k\)-abhängiges Feinspektrum.
\end{itemize}

Im Bamberg-Sinn bedeutet dies: \(Q_0\) ist die Urform der Viererordnung, während \(Q_k\) für \(k\ge 1\) strukturtreue Dehnungsnachfahren derselben Rahmenarchitektur sind.

\subsection{Kompakte Forschungsformel}

Die Primzahlvierlinge
\[
Q_0(p)=(p,p+2,p+6,p+8)
\]
bilden die primale Grundform einer unendlichen Dehnungsfamilie
\[
Q_k(p)=(p,p+2,p+6+12k,p+8+12k).
\]
Jede Stufe erhält die Zwillingsflügelstruktur und die EABC-Klassenarchitektur. Der Dehnungsrang \(k\) steuert den metrischen Kreuzrest quadratisch, die Walter-Besetzungsasymmetrie linear und die Morley-Feinverzerrung in untergeordneter Weise. Damit entsteht eine hierarchische Familie strukturtreuer, aber zunehmend gedehnter Viererordnungen, deren globale Rahmenstruktur invariant bleibt, während ihre endlichen Korrekturterme systematisch mit dem Dehnungsrang anwachsen.
